\documentclass[tikz,border=4pt]{standalone}
\usepackage{ctex}
\begin{document}
\begin{tikzpicture}[
    x={(1cm,0cm)}, y={(0,1cm)}, z={(-0.5cm,-0.4cm)},
    line cap=round, line join=round,
    every label/.style={font=\small}
]

% 共用顶点函数：在每个 scope 内重画立方体
\newcommand{\drawcube}{
  \coordinate (A) at (0,0,0);
  \coordinate (B) at (2,0,0);
  \coordinate (C) at (2,0,2);
  \coordinate (D) at (0,0,2);
  \coordinate (Ap) at (0,2,0);
  \coordinate (Bp) at (2,2,0);
  \coordinate (Cp) at (2,2,2);
  \coordinate (Dp) at (0,2,2);
  % 可见棱
  \draw[gray,thick] (A) -- (B) -- (C) -- (D) -- cycle;
  \draw[gray,thick] (Ap) -- (Bp) -- (Cp) -- (Dp) -- cycle;
  \draw[gray,thick] (A) -- (Ap);
  \draw[gray,thick] (B) -- (Bp);
  \draw[gray,thick] (C) -- (Cp);
  \draw[gray,thick] (D) -- (Dp);
}

% ============ 左：三角形截面（过 A、B、Dp 三顶点）============
\begin{scope}[xshift=0cm]
  \drawcube
  % 截面
  \fill[red!20,opacity=0.6] (A) -- (B) -- (Dp) -- cycle;
  \draw[red,thick] (A) -- (B) -- (Dp) -- cycle;
  \node[red,below] at (A) {$A$};
  \node[red,below right] at (B) {$B$};
  \node[red,above] at (Dp) {$D'$};
  \node[font=\bfseries] at (1,-1.2,1) {① 三角形截面};
\end{scope}

% ============ 中：矩形截面（中横截）============
\begin{scope}[xshift=5cm]
  \drawcube
  \coordinate (M1) at (1,0,0);
  \coordinate (M2) at (1,0,2);
  \coordinate (M3) at (1,2,2);
  \coordinate (M4) at (1,2,0);
  \fill[blue!20,opacity=0.6] (M1) -- (M2) -- (M3) -- (M4) -- cycle;
  \draw[blue,thick] (M1) -- (M2) -- (M3) -- (M4) -- cycle;
  \node[font=\bfseries] at (1,-1.2,1) {② 矩形截面};
\end{scope}

% ============ 右：六边形截面（过 6 棱中点）============
\begin{scope}[xshift=10cm]
  \drawcube
  \coordinate (P1) at (1,0,0);   % AB 中
  \coordinate (P2) at (2,1,0);   % BBp 中
  \coordinate (P3) at (2,2,1);   % BpCp 中
  \coordinate (P4) at (1,2,2);   % CpDp 中
  \coordinate (P5) at (0,1,2);   % DDp 中
  \coordinate (P6) at (0,0,1);   % AD 中
  \fill[green!20,opacity=0.6] (P1) -- (P2) -- (P3) -- (P4) -- (P5) -- (P6) -- cycle;
  \draw[green!50!black,thick] (P1) -- (P2) -- (P3) -- (P4) -- (P5) -- (P6) -- cycle;
  \node[font=\bfseries] at (1,-1.2,1) {③ 六边形截面};
\end{scope}

\end{tikzpicture}
\end{document}
